\histitem{Discovery of the continuous spectrum}

In the Note on topological vector spaces (cf. ÉHM, pp. 262–263), we described how I. Fredholm, around 1900 — continuing the work of C. Neumann, V. Volterra, and H. Poincaré — came to consider the equation

\[
u (x) - \lambda \int_ {\mathrm{I}} \mathrm{K} (x, y) u (y) d y = f (x)\tag{1}
\]

with unknown function \(u\colon I\to\mathbf{C}\) , in the case where the interval \(I\subset\mathbf{R}\) is compact and where the kernel \(K\) is continuous on \(I\times I\) . A skillful use of "infinite determinants," based on ideas of Poincaré and H. von Koch, allowed him to obtain families of solutions whose dependence on the complex variable \(\lambda\) is meromorphic, then to prove the famous “alternative” regarding the existence of solutions (cf.~III, p.~81, th. 5). But Fredholm immediately specializes his results to the case \(\lambda = -1\) without exploiting the fact that equation (1) is an eigenvalue problem \(^{(1)}\) .

When the kernel K is symmetric, we also saw how Hilbert, in his first memoir on integral equations [31], recognized the kinship between the Fredholm problem and the search for the “principal axes” of a quadratic form. Starting from this reduction for finite “sections” (discretizations) of the kernel K, he obtained by passage to the limit the formula

\[
\int_ {\mathrm{I}} \mathrm{K} (s, t) x (s) x (t) d t = \sum_ {n = 1} ^ {\infty} \frac {1}{\lambda_ {n}} \int_ {\mathrm{I}} \varphi_ {n} (s) x (s) d s,\tag{2}
\]

where the \(\lambda_{n}\) are the eigenvalues of the kernel K, the \(\varphi_{n}\) form an orthonormal system of associated eigenvectors \(^{(2)}\) and where the equality holds as soon as x is square-integrable on I (cf.~ÉHM, p.~264).

The principal obstacles facing Hilbert concerned the passage from finite sums to the integrals of formula (2). He deduced from this formula the existence of the eigenvalues — however, he gave for the latter a variational characterization independent of the passage to the limit, inspired by the classical method for determining eigenvalues in finite dimension (cf.~no. 4 of IV, p.~153). E. Schmidt would show in 1905, drawing on the work of H. Schwarz, how to obtain (2) without going through the “infinite determinants” on which the methods of Fredholm and Hilbert depended (cf.~ÉHM, p.~265).

But already in 1906, Hilbert [32] turns to the more general problem of reducing a quadratic form

\[
\mathrm{B} (x, x) = \sum_ {p, q = 0} ^ {\infty} b _ {p q} x _ {p} x _ {q}, \qquad x = (x _ {p}) _ {p \in \mathbf {N}}
\]

on \(E = \ell_{\mathbf{C}}^{2}(\mathbf{N})\) . He certainly has in mind that by passing to coefficients in an orthonormal basis of E, problem (1) reduces to the search for solutions of the “infinite linear system”

\[
x _ {p} + \sum_ {q = 0} ^ {\infty} b _ {p q} x _ {q} = f _ {p} \quad (p \in \mathbf {N}).
\]

He notes, however, that for the Fredholm problem one deals with very particular quadratic forms, for which \(\mathrm{B}(x_{n}, x_{n})\) tends to \(\mathrm{B}(x, x)\) as soon as \(x_{n}\) converges weakly to x; he calls them “completely continuous” (vollstetig).

Hilbert insists on the essential difference that distinguishes them from quadratic forms bounded on E, and decides to undertake the more general reduction of the latter. Again, his method is to start from the reduction of the "finite sections."

\[
(x _ {0}, \ldots , x _ {n}) \mapsto \sum_ {p, q = 0} ^ {n} b _ {p q} x _ {p} x _ {q}
\]

of the form B, and then to let n tend to infinity. For this passage to the limit, he takes advantage of the ideas on integration introduced by T. Stieltjes in his work on continued fractions [71]. He shows that every bounded quadratic form can, after an orthogonal transformation of the variables, be put in the form

\[
\sum_ {p \in {\bf N}} \frac {1}{\lambda_ {p}} x _ {p} ^ {2} + \int_ {s \in {\bf R}} \frac {1}{s} d \sigma (s, x).\tag{3}
\]

In this formula, the \(\lambda_{p}\) are the eigenvalues of B and \(x\mapsto\sigma(s,x)\) is (for a fixed \(s\in\mathbf{R}\)) a positive definite separating quadratic form on E, while \(s\mapsto\sigma(s,x)\) is \(^{(3)}\) (for a fixed \(x\in\mathrm{E}\)) a continuous increasing function that tends to 0 at \(-\infty\) and toward \(\|x\|^2\) at \(+\infty\) (cf. ÉHM, p. 284, for Stieltjes notation).

Let S denote the set of points of \(\overline{R}\) that do not admit a neighborhood in which \(s \mapsto \sigma(s, x)\) remains constant for all x; then the integral in (3) can naturally be taken over S. Hilbert calls the set S the “continuous spectrum” (Streckenspektrum), the set of eigenvalues of B the “point spectrum” (Punktspektrum), and their union the “spectrum” (Spektrum). This term, coming from optics, had been used in 1897 by W. Wirtinger [91] in the study of differential equations with periodic coefficients.

Hilbert would quickly draw from his “spectral” methods a profound renewal of several problems of classical analysis, pertaining to the study of ordinary differential equations (in particular of Sturm–Liouville type), partial differential equations, or functions of a complex variable. We will not report here on the profusion of results that followed in the first quarter of the twentieth century and refer to the synthesis by E. Hellinger and O. Toeplitz [29]. We shall nevertheless mention some of these works on integral equations, in which the germs of later developments of the abstract theory are to be found.

For Hilbert and his students, faithful to a German tradition in linear algebra, the model remains the principal axes theorem; thus it is always a matter of studying the spectra of bilinear or quadratic forms, in particular on the “Hilbert” space \(E = \ell_{\mathbf{C}}^{2}(\mathbf{N})\) . The fact that one can associate any continuous bilinear form B on E with the continuous endomorphism A of E characterized by \(\langle Ax, y \rangle = B(x, y)\) was known to the Hilbert school, particularly after the works of E. Schmidt, M. Fréchet, and F. Riesz in 1907–1908. But it was F. Riesz who showed, in an admirable book [59] from 1913, the interest of placing endomorphisms at the forefront in this context.

Riesz gives the modern definition of the spectrum of an element A of \(\mathcal{L}(\mathrm{E})\) [59, p.~139], notes that this is a closed, bounded set, and shows that the resolvent \(\lambda \mapsto (\lambda 1_{\mathrm{E}} - \mathrm{A})^{-1}\) is holomorphic on the complement of the spectrum of A. Above all, he assigns a central role to the algebra structure of \(\mathcal{L}(\mathrm{E})\): Drawing on the work of Volterra, he developed a first version of the functional calculus for symmetric endomorphisms, which allowed him to recover simply the reduction results of Hilbert and Schmidt. If A is a symmetric element of \(\mathcal{L}(\mathrm{E})\) , Riesz shows how to define \(f(\mathrm{A})\) for every function \(f\) semicontinuous (from below or from above) on \(\mathbf{R}\) in \(\mathbf{R}\) , and notes that \(f \mapsto f(\mathrm{A})\) is what we today call an algebra morphism [59, p.~129--130]. Applying this idea to the characteristic function of an interval \([\xi, +\infty[\) , one obtains for every \(\xi\) a symmetric endomorphism \(\mathrm{A}_{\xi}\) of E; using the Stieltjes integral, Riesz then proves a version of Hilbert's relation (3):

\[
\langle \mathrm{A} x, y \rangle = \int_ {\mathbf {R}} \xi d \langle \mathrm{A} _ {\xi} x, y \rangle
\]

for \(x, y \in \mathrm{E}\) , an equality that he even writes

\[
\mathrm{A} = \int_ {\mathbf {R}} \xi d \mathrm{A} _ {\xi}.
\]

The spectrum of A is then the set of values \(\xi_{0}\) such that \(A_{\xi}\) does not remain constant in any neighborhood of \(\xi_{0}\) and the point spectrum that of the points of discontinuity of \(\xi \mapsto A_{\xi}\) .

Shortly before concluding with the applications of the theory, Riesz notes [59, p.~146] that its results take a particularly simple form when the endomorphism A arises from one of Hilbert's "completely continuous" bilinear forms: the spectrum reduces to the set of eigenvalues (to which 0 must possibly be added), the eigenvalues are arranged in a sequence tending to 0, and each nonzero eigenvalue has finite multiplicity (III, p.~90, prop. 5).

